\chapter{Bounded operators on Hilbert space}\label{ch:hilbert}

This appendix collects the properties of bounded operators on complex Hilbert
spaces which are used throughout the text: elementary facts about the algebra of
bounded operators, compact operators, Fredholm operators and their index, and the
Hilbert--Schmidt and trace-class ideals. Proofs are short and standard. Hilbert
spaces $H$, $H_1$, $H_2,\dots$ are complex; where orthonormal bases are used they are
assumed separable, so that bases are countable.

\section{Bounded operators}\label{hilbert:sec:bounded}

\begin{proposition}\label{hilbert:prop:bounded}
A linear map $A:H_1\to H_2$ is continuous if and only if it is bounded, i.e.\
$\|A\|=\sup_{\|x\|\le1}\|Ax\|<\infty$. The space $\Lin(H_1,H_2)$ of bounded operators is a
Banach space with this norm, and $\|AB\|\le\|A\|\,\|B\|$ whenever the composite is
defined. In particular $\Lin(H)=\Lin(H,H)$ is a Banach algebra.
\end{proposition}
\begin{proof}
If $A$ is continuous at $0$ there is $\delta>0$ with $\|Ax\|\le1$ for $\|x\|\le\delta$, so
$\|A\|\le\delta^{-1}$; conversely $\|Ax-Ay\|\le\|A\|\|x-y\|$. If $(A_k)$ is Cauchy in norm then
$A_kx$ is Cauchy for each $x$, and the limit $Ax$ is linear with $\|A-A_k\|\to0$. Finally
$\|ABx\|\le\|A\|\|Bx\|\le\|A\|\|B\|\|x\|$.
\end{proof}

By the Riesz representation theorem each $A\in\Lin(H_1,H_2)$ has an adjoint
$A^*\in\Lin(H_2,H_1)$, $\inner{Ax}y=\inner x{A^*y}$, with $(AB)^*=B^*A^*$, $A^{**}=A$ and
$\|A^*\|=\|A\|$ (since $\|A\|=\sup_{\|x\|,\|y\|\le1}|\inner{Ax}y|$). We also use the open
mapping theorem: a bounded bijection between Banach spaces has bounded inverse.

\begin{proposition}[$C^*$-identity]\label{hilbert:prop:Cstar}
For $A\in\Lin(H_1,H_2)$, $\|A^*A\|=\|AA^*\|=\|A\|^2$.
\end{proposition}
\begin{proof}
$\|Ax\|^2=\inner{A^*Ax}x\le\|A^*A\|\|x\|^2$, so $\|A\|^2\le\|A^*A\|\le\|A^*\|\|A\|=\|A\|^2$. Apply this
to $A^*$ for the other identity.
\end{proof}

\begin{proposition}[Neumann series]\label{hilbert:prop:neumann}
If $A\in\Lin(H)$ and $\|A\|<1$ then $\Id-A$ is invertible, with
$(\Id-A)^{-1}=\sum_{k\ge0}A^k$ and $\|(\Id-A)^{-1}\|\le(1-\|A\|)^{-1}$. Consequently:
\begin{enumerate}[label=\textup{(\alph*)}]
\item $A-\lambda$ is invertible for $|\lambda|>\|A\|$, so the spectrum of $A$ lies in
$\{|\lambda|\le\|A\|\}$;
\item the invertible elements of $\Lin(H_1,H_2)$ form an open set: if $A$ is invertible and
$\|E\|<\|A^{-1}\|^{-1}$ then $A+E=A(\Id+A^{-1}E)$ is invertible; and inversion is continuous.
\end{enumerate}
\end{proposition}
\begin{proof}
The series converges absolutely in the Banach space $\Lin(H)$ since $\|A^k\|\le\|A\|^k$, and
$(\Id-A)\sum_{k<N}A^k=\sum_{k<N}A^k(\Id-A)=\Id-A^N\to\Id$. For (a) write
$A-\lambda=-\lambda(\Id-A/\lambda)$. In (b) the inverse is
$(A+E)^{-1}=\sum_k(-A^{-1}E)^kA^{-1}$, which depends continuously on $E$.
\end{proof}

The condition in (a) cannot be weakened to $|\lambda|\ge\|A\|$: for $A=\Id$, $\lambda=1$.

\begin{lemma}\label{hilbert:lem:selfadjoint-norm}
If $T\in\Lin(H)$ is self-adjoint then $\|T\|=\sup_{\|x\|=1}|\inner{Tx}x|$.
\end{lemma}
\begin{proof}
Let $M$ be the supremum; clearly $M\le\|T\|$. By polarization, for $\|x\|=\|y\|=1$,
$4\operatorname{Re}\inner{Tx}y=\inner{T(x+y)}{x+y}-\inner{T(x-y)}{x-y}\le M(\|x+y\|^2+\|x-y\|^2)=4M$.
Taking $y=Tx/\|Tx\|$ gives $\|Tx\|\le M$.
\end{proof}

\section{Compact operators}\label{hilbert:sec:compact}

\begin{definition}\label{hilbert:def:compact}
An operator $K\in\Lin(H_1,H_2)$ is \emph{compact} if the closure of the image of the
unit ball is compact; equivalently, if every bounded sequence $(x_n)$ has a
subsequence along which $Kx_n$ converges. The compact operators form a subspace
$\Kpt(H_1,H_2)$, and $\Kpt(H)=\Kpt(H,H)$.
\end{definition}

\begin{proposition}\label{hilbert:prop:compact}
\begin{enumerate}[label=\textup{(\alph*)}]
\item If $K$ is compact and $A$, $B$ are bounded, then $AKB$ is compact (when defined);
$\Kpt(H)$ is a two-sided ideal in $\Lin(H)$.
\item $K$ is compact if and only if $x_n\rightharpoonup x$ weakly implies $Kx_n\to Kx$ in norm.
\item $\Kpt(H_1,H_2)$ is the norm closure of the finite-rank operators.
\item $K$ is compact if and only if $K^*$ is compact, if and only if $K^*K$ is compact.
\item \textup{(Riesz)} If the closed unit ball of $H$ is compact, $\dim H<\infty$. In particular
if $K$ is compact, $\Nul(\Id-K)$ is finite-dimensional.
\end{enumerate}
\end{proposition}
\begin{proof}
(a) Bounded operators map bounded sets to bounded sets and convergent sequences to
convergent sequences.

(b) Suppose $K$ compact and $x_n\rightharpoonup x$. Then $(x_n)$ is bounded and $Kx_n\rightharpoonup Kx$,
since $\inner{Kx_n}z=\inner{x_n}{K^*z}$. Every subsequence of $(Kx_n)$ has a norm-convergent
subsequence, whose limit must be $Kx$; hence $Kx_n\to Kx$. Conversely, every bounded sequence
has a weakly convergent subsequence, which $K$ maps to a convergent sequence.

(c) Finite-rank operators are compact, and a norm limit $K$ of compact operators $K_j$ is
compact: given $\eps>0$ choose $j$ with $\|K-K_j\|<\eps$; then $K(B)$ lies within $\eps$ of the
totally bounded set $K_j(B)$, so is totally bounded. Conversely, let $K$ be compact and $\eps>0$.
Cover the compact set $\overline{K(B)}$ by finitely many balls $B(y_i,\eps)$, and let $P$ be the
orthogonal projection onto the span of the $y_i$. For $\|x\|\le1$ choose $i$ with
$\|Kx-y_i\|<\eps$; then $\|Kx-PKx\|\le\|Kx-y_i\|+\|P(y_i-Kx)\|<2\eps$. So $\|K-PK\|\le2\eps$, and
$PK$ has finite rank.

(d) The adjoint of a finite-rank operator has finite rank and $\|K^*-F^*\|=\|K-F\|$, so the
first equivalence follows from (c). If $K$ is compact so is $K^*K$, by (a). Conversely,
suppose $K^*K$ is compact and $x_n\rightharpoonup x$, $\|x_n\|\le M$. Then
$\|K(x_n-x)\|^2=\inner{K^*K(x_n-x)}{x_n-x}\le2M\|K^*K(x_n-x)\|\to0$ by (b), so $K$ is compact by (b).

(e) If $\dim H=\infty$ an orthonormal sequence $(e_n)$ has $\|e_n-e_m\|=\sqrt2$ for $n\ne m$, so
has no convergent subsequence. On $\Nul(\Id-K)$, $x=Kx$, so the unit ball of this closed
subspace is compact.
\end{proof}

\begin{theorem}[Spectral theorem]\label{hilbert:thm:spectral}
Let $T\in\Kpt(H)$ be self-adjoint. Then there is an orthonormal sequence $(e_k)$ (finite
or infinite) of eigenvectors, $Te_k=\lambda_ke_k$, with real eigenvalues
$|\lambda_1|\ge|\lambda_2|\ge\dots>0$ tending to zero if infinite, such that
$H=\Nul(T)\oplus\overline{\operatorname{span}}\{e_k\}$ and $Tx=\sum_k\lambda_k\inner x{e_k}e_k$.
Each non-zero eigenvalue has finite multiplicity.
\end{theorem}
\begin{proof}
If $T\ne0$, by \cref{hilbert:lem:selfadjoint-norm} there are unit vectors $x_n$ with
$\inner{Tx_n}{x_n}\to\lambda$, $|\lambda|=\|T\|$, $\lambda$ real. Then
$\|Tx_n-\lambda x_n\|^2=\|Tx_n\|^2-2\lambda\inner{Tx_n}{x_n}+\lambda^2\le2\lambda^2-2\lambda\inner{Tx_n}{x_n}\to0$.
Passing to a subsequence, $Tx_n\to y$, so $x_n\to y/\lambda$, and $e=y/\|y\|$ satisfies $Te=\lambda e$.
The orthocomplement of $e$ is invariant under $T$; repeating the argument there gives the
sequence $(e_k)$ with $|\lambda_k|$ non-increasing, the restriction of $T$ to the orthocomplement
of $e_1,\dots,e_k$ having norm $|\lambda_{k+1}|$. If infinitely many $|\lambda_k|\ge\eps>0$, then
$\|Te_k-Te_j\|^2=\lambda_k^2+\lambda_j^2\ge2\eps^2$, contradicting compactness; so $\lambda_k\to0$, and each
eigenvalue occurs finitely often. On the orthocomplement $W$ of all the $e_k$, which is
invariant, $\|T|_W\|\le|\lambda_k|$ for all $k$ (or the process would have stopped), so
$T|_W=0$ and $W\subset\Nul(T)$; the reverse inclusion holds since
$\inner x{e_k}=\lambda_k^{-1}\inner{x}{Te_k}=\lambda_k^{-1}\inner{Tx}{e_k}$.
\end{proof}

\section{Fredholm operators}\label{hilbert:sec:fredholm}

\begin{definition}\label{hilbert:def:fredholm}
An operator $A\in\Lin(H_1,H_2)$ is \emph{Fredholm} if $\Nul(A)$ is finite-dimensional,
$\Ran(A)$ is closed, and $\Ran(A)^\perp$ is finite-dimensional. Its \emph{index} is
$\ind(A)=\dim\Nul(A)-\dim\Ran(A)^\perp=\dim\Nul(A)-\dim\Nul(A^*)$.
\end{definition}

Here $\Ran(A)^\perp=\Nul(A^*)$ for any bounded $A$, and if $\Ran(A)$ is closed it is the
orthocomplement of $\Nul(A^*)$. Thus $A$ is Fredholm if and only if $A^*$ is, and
$\ind(A^*)=-\ind(A)$.

\begin{discussion}
Fredholm operators are the operators that are ``invertible up to finite
dimensions''. Their null space and cokernel can jump under perturbation, but
only together, so that the difference of their dimensions, the index, is
stable. Atkinson's theorem below makes this practical: to prove that an
operator is Fredholm it suffices to invert it modulo compact operators, which
is exactly what a parametrix does on a compact manifold.
\end{discussion}

\begin{proposition}[Generalized inverse]\label{hilbert:prop:generalized-inverse}
If $A\in\Lin(H_1,H_2)$ is Fredholm there is a unique $G\in\Lin(H_2,H_1)$ with
\begin{equation*}
GA=\Id-\pi_N,\qquad AG=\Id-\pi_C,\qquad G\pi_C=0,
\end{equation*}
where $\pi_N$ and $\pi_C$ are the orthogonal projections onto $\Nul(A)$ and $\Nul(A^*)$. $G$ is
Fredholm with $\ind(G)=-\ind(A)$.
\end{proposition}
\begin{proof}
$A$ restricts to a bounded bijection $\Nul(A)^\perp\to\Ran(A)$ between Hilbert spaces, which has
bounded inverse $A_0^{-1}$ by the open mapping theorem; set $G=A_0^{-1}(\Id-\pi_C)$. The
identities are immediate and determine $G$. $\Nul(G)=\Nul(A^*)$ and $\Ran(G)=\Nul(A)^\perp$, so
$\ind(G)=\dim\Nul(A^*)-\dim\Nul(A)$.
\end{proof}

\begin{theorem}[Atkinson]\label{hilbert:thm:atkinson}
$A\in\Lin(H_1,H_2)$ is Fredholm if and only if there exists $B\in\Lin(H_2,H_1)$ such that
$BA-\Id$ and $AB-\Id$ are compact. One may then take $B$ such that both are of finite rank.
\end{theorem}
\begin{proof}
If $A$ is Fredholm, the generalized inverse is such a $B$, with finite-rank errors.
Conversely suppose $BA=\Id-K$ and $AB=\Id-K'$ with $K,K'$ compact. Then
$\Nul(A)\subset\Nul(\Id-K)$, which is finite-dimensional by \cref{hilbert:prop:compact}(e).
Let $L=\Id-K'$. On $\Nul(L)^\perp$ there is $c>0$ with $\|Lx\|\ge c\|x\|$: otherwise there are unit
vectors $x_n\in\Nul(L)^\perp$ with $Lx_n\to0$; passing to a subsequence $K'x_n\to y$, so
$x_n=Lx_n+K'x_n\to y$, with $\|y\|=1$, $y\in\Nul(L)^\perp$ and $Ly=0$, a contradiction. Hence
$\Ran(L)$ is closed, and $\Ran(L)^\perp=\Nul(\Id-K'^*)$ is finite-dimensional. Since
$\Ran(A)\supset\Ran(AB)=\Ran(L)$, $\Ran(A)$ is the sum of the closed subspace $\Ran(L)$ and a
finite-dimensional space, hence closed, and $\Ran(A)^\perp\subset\Ran(L)^\perp$ is
finite-dimensional.
\end{proof}

\begin{lemma}\label{hilbert:lem:finite-rank-index}
If $F\in\Lin(H)$ has finite rank then $\Id-F$ is Fredholm of index zero.
\end{lemma}
\begin{proof}
Let $V=\Ran(F)+\Ran(F^*)$, which is finite-dimensional. Then $F$ maps $V$ into itself and
vanishes on $V^\perp\subset\Nul(F)$, so $\Id-F=(\Id_V-F|_V)\oplus\Id_{V^\perp}$, and the index is that
of an endomorphism of a finite-dimensional space, which is zero.
\end{proof}

\begin{proposition}\label{hilbert:prop:index-product}
If $B\in\Lin(H_1,H_2)$ and $A\in\Lin(H_2,H_3)$ are Fredholm, then $AB$ is Fredholm and
$\ind(AB)=\ind(A)+\ind(B)$.
\end{proposition}
\begin{proof}
If $B'$, $A'$ are parametrices as in \cref{hilbert:thm:atkinson} then $B'A'$ is one for $AB$.
Consider the sequence of finite-dimensional spaces and natural maps
\begin{equation*}
0\to\Nul B\hookrightarrow\Nul AB\xrightarrow{B}\Nul A\to H_2/\Ran B\xrightarrow{A}H_3/\Ran AB\to H_3/\Ran A\to0 ,
\end{equation*}
where the third map is inclusion followed by the quotient map and the last is the
quotient map ($\Ran AB\subset\Ran A$). It is exact: the kernel of $B$ on $\Nul AB$ is $\Nul B$; the
image of $\Nul AB$ under $B$ is $\Nul A\cap\Ran B$, the kernel of the third map; the kernel of
the map induced by $A$ consists of the classes of $v$ with $Av\in\Ran AB$, i.e.\
$v\in\Nul A+\Ran B$; and the kernel of the last map is $\Ran A/\Ran AB$. The alternating sum of the
dimensions in an exact sequence vanishes, and $\dim H_2/\Ran B=\dim\Nul B^*$ etc.\ since the
ranges are closed; this gives the formula.
\end{proof}

\begin{question}
Why must the range be closed in the definition? Is it not enough that
$\Nul(A)$ and $\Ran(A)^\perp$ are finite-dimensional?
\end{question}
\begin{answer}
No. The diagonal operator $Ae_k=k^{-1}e_k$ on a space with orthonormal basis
$(e_k)$ is injective with dense range, so both spaces are zero; but it is
compact, its range is not closed, and it has no bounded inverse on its range.
Being compact, it is certainly not invertible modulo compact operators, so
Atkinson's theorem fails for it. Closedness of the range is what the
parametrix characterization and the stability of the index rely on.
\end{answer}

\begin{theorem}[Stability of the index]\label{hilbert:thm:stability}
Let $A\in\Lin(H_1,H_2)$ be Fredholm.
\begin{enumerate}[label=\textup{(\alph*)}]
\item If $G$ is a generalized inverse of $A$ and $\|E\|<\|G\|^{-1}$, then $A+E$ is Fredholm and
$\ind(A+E)=\ind(A)$. Thus the Fredholm operators form an open subset of $\Lin(H_1,H_2)$ on
which the index is locally constant, hence constant along continuous paths.
\item If $K\in\Kpt(H_1,H_2)$ then $A+K$ is Fredholm and $\ind(A+K)=\ind(A)$.
\end{enumerate}
\end{theorem}
\begin{proof}
(a) Since $\|GE\|<1$ and $\|EG\|<1$, $\Id+GE$ and $\Id+EG$ are invertible, and
\begin{gather*}
(\Id+GE)^{-1}G(A+E)=\Id-(\Id+GE)^{-1}\pi_N,\\ (A+E)G(\Id+EG)^{-1}=\Id-\pi_C(\Id+EG)^{-1},
\end{gather*}
using $GA=\Id-\pi_N$ and $AG=\Id-\pi_C$. By \cref{hilbert:thm:atkinson}, $A+E$ is Fredholm. The
first identity, \cref{hilbert:lem:finite-rank-index} and \cref{hilbert:prop:index-product}
(invertible operators having index zero) give $\ind(G)+\ind(A+E)=0$, and
$\ind(G)=-\ind(A)$.

(b) A parametrix for $A$ is a parametrix for $A+tK$ modulo compact operators, so $A+tK$ is
Fredholm for $t\in[0,1]$, and by (a) the index is locally constant, hence constant, in $t$.
\end{proof}

\begin{theorem}[Riesz theory]\label{hilbert:thm:riesz}
Let $K\in\Kpt(H)$ and $\mu\in\C\setminus\{0\}$. Then $\mu-K$ is Fredholm of index zero, so $\mu\in\Spec(K)$
if and only if $\mu$ is an eigenvalue. For each $\eps>0$ there are only finitely many
eigenvalues with $|\mu|\ge\eps$; thus $\Spec(K)\setminus\{0\}$ is countable and discrete in
$\C\setminus\{0\}$. For each eigenvalue $\mu\ne0$ the generalized eigenspace
$\bigcup_N\Nul((\mu-K)^N)$ is finite-dimensional.
\end{theorem}
\begin{proof}
$\mu-K=\mu(\Id-K/\mu)$ is Fredholm of index zero by \cref{hilbert:thm:stability}(b); if it is
injective it is therefore surjective, hence invertible. Suppose there were infinitely
many distinct eigenvalues $\mu_k$ with $|\mu_k|\ge\eps$ and eigenvectors $x_k$; these are linearly
independent. Let $M_k=\operatorname{span}\{x_1,\dots,x_k\}$ and choose unit vectors $y_k\in M_k\ominus M_{k-1}$.
Since $(K-\mu_k)M_k\subset M_{k-1}$, for $j<k$
\begin{equation*}
Ky_k-Ky_j=\mu_ky_k+\big((K-\mu_k)y_k-Ky_j\big),\quad (K-\mu_k)y_k-Ky_j\in M_{k-1},
\end{equation*}
so $\|Ky_k-Ky_j\|\ge|\mu_k|\ge\eps$, contradicting compactness. Finally the spaces
$N_j=\Nul((\mu-K)^j)$ increase; if $N_{j-1}\ne N_j$ for all $j$, choose unit $u_j\in N_j\ominus N_{j-1}$.
For $i<j$, $Ku_j-Ku_i=\mu u_j-\big((\mu-K)u_j+Ku_i\big)$ with the bracket in $N_{j-1}$, so
$\|Ku_j-Ku_i\|\ge|\mu|$, again a contradiction. So $N_j=N_{j_0}$ for $j\ge j_0$, and
$N_{j_0}=\Nul((\mu-K)^{j_0})$ is finite-dimensional because $(\mu-K)^{j_0}=\mu^{j_0}(\Id-K')$ with $K'$
compact.
\end{proof}

\section{Hilbert--Schmidt and trace-class operators}\label{hilbert:sec:trace}

\begin{definition}\label{hilbert:def:HS}
An operator $A\in\Lin(H_1,H_2)$ is \emph{Hilbert--Schmidt} if
$\|A\|_2^2=\sum_j\|Ae_j\|^2<\infty$ for an orthonormal basis $(e_j)$ of $H_1$.
\end{definition}

For orthonormal bases $(e_j)$ of $H_1$ and $(f_k)$ of $H_2$,
$\sum_j\|Ae_j\|^2=\sum_{j,k}|\inner{Ae_j}{f_k}|^2=\sum_k\|A^*f_k\|^2$, so $\|A\|_2$ does not depend on the
basis and $\|A^*\|_2=\|A\|_2$. The Hilbert--Schmidt operators form a Hilbert space with inner
product $\langle A,B\rangle_2=\sum_j\inner{Ae_j}{Be_j}$, also independent of the basis, and:
\begin{itemize}
\item $\|A\|\le\|A\|_2$ (expand $x$ in a basis and use Cauchy--Schwarz);
\item $\|SAT\|_2\le\|S\|\,\|A\|_2\,\|T\|$ for bounded $S,T$;
\item Hilbert--Schmidt operators are compact: if $P_N$ is the projection onto
$\operatorname{span}\{e_1,\dots,e_N\}$, $\|A-AP_N\|_2^2=\sum_{j>N}\|Ae_j\|^2\to0$.
\end{itemize}

\begin{proposition}\label{hilbert:prop:HS-kernel}
Let $(M_i,\mu_i)$, $i=1,2$, be $\sigma$-finite measure spaces with $L^2(M_i)$ separable. An
operator $A:L^2(M_1)\to L^2(M_2)$ is Hilbert--Schmidt if and only if it is given by a kernel
$k\in L^2(M_2\times M_1)$, $Au(x)=\int k(x,y)u(y)\,d\mu_1(y)$, and then $\|A\|_2=\|k\|_{L^2}$. The same
holds for vector-valued functions.
\end{proposition}
\begin{proof}
If $(e_j)$ and $(f_k)$ are orthonormal bases of $L^2(M_1)$ and $L^2(M_2)$, the products
$f_k(x)\overline{e_j(y)}$ form an orthonormal basis of $L^2(M_2\times M_1)$. If $A$ has kernel
$k\in L^2$, then $\inner{Ae_j}{f_k}=\langle k,f_k\otimes\bar e_j\rangle$, so
$\|A\|_2^2=\sum_{j,k}|\inner{Ae_j}{f_k}|^2=\|k\|_{L^2}^2$. Conversely, if $A$ is Hilbert--Schmidt then
$k=\sum_{j,k}\inner{Ae_j}{f_k}f_k\otimes\bar e_j\in L^2$ is a kernel for $A$, as one checks on the $e_j$.
\end{proof}

\begin{definition}\label{hilbert:def:trace-class}
An operator $A\in\Lin(H)$ is of \emph{trace class} if it is a finite sum
$A=\sum_{k}B_kC_k$ of products of Hilbert--Schmidt operators $C_k:H\to H_k'$, $B_k:H_k'\to H$. The
trace norm $\|A\|_1$ is the infimum of $\sum_k\|B_k\|_2\|C_k\|_2$ over such representations.
\end{definition}

The trace-class operators form a two-sided ideal in $\Lin(H)$, contained in the
compact operators, with $\|SAT\|_1\le\|S\|\,\|A\|_1\|T\|$ and $\|A\|\le\|A\|_1$; $\|\cdot\|_1$ is a norm,
and the ideal is complete with respect to it. Every finite-rank operator is of trace class.
(Using the polar decomposition $A=U|A|$ one can show that
$\|A\|_1=\sum_j\inner{|A|e_j}{e_j}$, the sum of the singular values of $A$.)

\begin{proposition}\label{hilbert:prop:trace}
If $A\in\Lin(H)$ is of trace class then $\Tr(A)=\sum_j\inner{Ae_j}{e_j}$ converges absolutely, is
independent of the orthonormal basis $(e_j)$, and $|\Tr(A)|\le\|A\|_1$. Moreover:
\begin{enumerate}[label=\textup{(\alph*)}]
\item $\Tr(XY)=\Tr(YX)$ if $X\in\Lin(H',H)$ and $Y\in\Lin(H,H')$ are Hilbert--Schmidt;
\item $\Tr(ST)=\Tr(TS)$ if $S$ is of trace class and $T$ is bounded;
\item for finite-rank $A$, $\Tr(A)$ is the trace of the matrix of $A$ restricted to any
finite-dimensional space containing $\Ran(A)$ and $\Ran(A^*)$;
\item if $A=BC$ on $L^2(M)$, with $B$, $C$ Hilbert--Schmidt with kernels $b$, $c$, then
$\Tr(A)=\int\!\!\int b(x,y)c(y,x)\,d\mu(y)\,d\mu(x)$.
\end{enumerate}
\end{proposition}
\begin{proof}
For $A=BC$, $\sum_j\inner{BCe_j}{e_j}=\sum_j\inner{Ce_j}{B^*e_j}=\langle C,B^*\rangle_2$, which converges
absolutely by Cauchy--Schwarz with modulus at most $\|B\|_2\|C\|_2$, and is independent of the basis.
Summing over the terms of a representation gives the first statements.

(a) For orthonormal bases $(e_j)$ of $H$ and $(f_k)$ of $H'$,
\begin{equation*}
\Tr(XY)=\sum_j\inner{Ye_j}{X^*e_j}=\sum_{j,k}\inner{Ye_j}{f_k}\inner{f_k}{X^*e_j}
=\sum_{j,k}\inner{Xf_k}{e_j}\inner{e_j}{Y^*f_k}=\Tr(YX),
\end{equation*}
the double sum converging absolutely by Cauchy--Schwarz.

(b) It suffices to take $S=BC$; by (a), $\Tr(BCT)=\Tr(CTB)=\Tr(TBC)$.

(c) Take a basis containing an orthonormal basis of that space; $A$ vanishes on its
orthocomplement.

(d) $\Tr(BC)=\langle C,B^*\rangle_2$, and the kernel of $B^*$ is $\overline{b(y,x)}$, so by
\cref{hilbert:prop:HS-kernel} this is $\int\!\!\int c(y,x)b(x,y)$.
\end{proof}

\begin{remark}[Schatten classes]\label{hilbert:rem:schatten}
For a compact operator $A$ let $s_1\ge s_2\ge\dots\ge0$ be its singular values, the eigenvalues
of $(A^*A)^{1/2}$. For $1\le p<\infty$ the Schatten class $\mathcal S_p$ consists of those $A$ with
$\|A\|_p=(\sum_js_j^p)^{1/p}<\infty$. These are two-sided ideals in $\Lin(H)$, Banach spaces under
$\|\cdot\|_p$, with $\mathcal S_1$ the trace class, $\mathcal S_2$ the Hilbert--Schmidt operators,
$\mathcal S_p\subset\mathcal S_q$ for $p\le q$, and H\"older's inequality
$\|AB\|_r\le\|A\|_p\|B\|_q$ for $\frac1r=\frac1p+\frac1q$. By Lidskii's theorem the trace of an operator
of trace class is the sum of its eigenvalues, counted with algebraic multiplicity.
\end{remark}
